% !TEX program = xelatex
% Самодостаточный конспект. Компилировать XeLaTeX/LuaLaTeX (tectonic 13-ryady-furye.tex).
\documentclass[11pt,a4paper]{article}
\newcommand{\coursename}{Математический анализ}
\newcommand{\rightheader}{§13. Ряды Фурье}
\newcommand{\doctitle}{Ряды Фурье}
% ==== Общая преамбула конспектов (собирается в каждый .tex как самодостаточная) ====
% Перед этим файлом должны быть определены: \documentclass и макросы
%   \coursename  — название курса (левый колонтитул, подзаголовок),
%   \rightheader — правый колонтитул (напр. «§2. Рациональные дроби»),
%   \doctitle    — заголовок раздела на титуле.
\usepackage{fontspec}
\setmainfont{cmunrm.otf}[
  BoldFont       = cmunbx.otf,
  ItalicFont     = cmunti.otf,
  BoldItalicFont = cmunbi.otf]
\usepackage{polyglossia}
\setdefaultlanguage{russian}
\setotherlanguage{english}

\usepackage{amsmath,amssymb,amsthm,mathtools}
\usepackage[a4paper,margin=2.2cm]{geometry}
\usepackage{enumitem}
\usepackage{graphicx}
\usepackage{array}
\usepackage{xcolor}
\definecolor{accent}{HTML}{1F6FEB}
\definecolor{rulecol}{HTML}{D0D7DE}
\usepackage[colorlinks=true,linkcolor=accent,urlcolor=accent,citecolor=accent]{hyperref}
\usepackage{titlesec}
\usepackage{fancyhdr}

% ----- Колонтитулы -----
\pagestyle{fancy}
\fancyhf{}
\fancyhead[L]{\small\itshape \coursename}
\fancyhead[R]{\small\itshape \rightheader}
\fancyfoot[C]{\thepage}
\renewcommand{\headrulewidth}{0.4pt}
\renewcommand{\headrule}{\hbox to\headwidth{\color{rulecol}\leaders\hrule height \headrulewidth\hfill}}

% ----- Заголовки -----
\titleformat{\section}{\large\bfseries\color{accent}}{\thesection.}{0.6em}{}
\titleformat{\subsection}{\bfseries}{\thesubsection.}{0.5em}{}

% ----- Теоремные окружения (стиль конспекта) -----
\theoremstyle{definition}
\newtheorem{definition}{Определение}[section]
\newtheorem{example}[definition]{Пример}
\newtheorem{remark}[definition]{Замечание}
\theoremstyle{plain}
\newtheorem{theorem}[definition]{Теорема}
\newtheorem{lemma}[definition]{Лемма}
\newtheorem{corollary}[definition]{Следствие}
\newtheorem{property}[definition]{Свойство}
\newtheorem{criterion}[definition]{Критерий}
\renewcommand{\proofname}{Доказательство}

% ----- Операторы и сокращения -----
\DeclareMathOperator{\tg}{tg}
\DeclareMathOperator{\ctg}{ctg}
\DeclareMathOperator{\arctg}{arctg}
\DeclareMathOperator{\arcctg}{arcctg}
\DeclareMathOperator{\sh}{sh}
\DeclareMathOperator{\ch}{ch}
\DeclareMathOperator{\Th}{th}
\DeclareMathOperator{\cth}{cth}
\DeclareMathOperator{\Const}{const}
\DeclareMathOperator{\sign}{sign}
\DeclareMathOperator{\rang}{rang}
\DeclareMathOperator{\diam}{diam}
\DeclareMathOperator{\Span}{span}
\DeclareMathOperator{\Ker}{Ker}
\DeclareMathOperator{\Img}{Im}
\DeclareMathOperator{\tr}{tr}
\newcommand{\R}{\mathbb{R}}
\newcommand{\C}{\mathbb{C}}
\newcommand{\N}{\mathbb{N}}
\newcommand{\Z}{\mathbb{Z}}
\newcommand{\Q}{\mathbb{Q}}
\newcommand{\dx}{\,dx}
\newcommand{\dt}{\,dt}
\newcommand{\du}{\,du}
\newcommand{\eps}{\varepsilon}
\renewcommand{\Re}{\operatorname{Re}}
\renewcommand{\Im}{\operatorname{Im}}
\renewcommand{\le}{\leqslant}
\renewcommand{\ge}{\geqslant}
\newcommand{\scal}[2]{\left( #1,\,#2\right)}
\DeclarePairedDelimiter{\abs}{\lvert}{\rvert}
\DeclarePairedDelimiter{\norm}{\lVert}{\rVert}

\begin{document}
\begin{center}
  {\LARGE\bfseries \doctitle}\\[4pt]
  {\large\itshape \coursename}
\end{center}
\vspace{6pt}

\section{Евклидово пространство функций. Пространство $L_2$}

Пусть $E$ — множество функций, заданных на отрезке $[a,b]$.

\begin{definition}
$E$ называется \emph{евклидовым} (предгильбертовым) пространством, если на нём задано
\emph{скалярное произведение} $(f,g)$ — билинейная, симметричная, положительно определённая форма.
В пространстве $L_2[a,b]$ интегрируемых с квадратом функций
\[
  (f,g)=\int_a^b f(x)g(x)\dx,\qquad
  \norm{f}_2=\sqrt{(f,f)}=\sqrt{\int_a^b f^2(x)\dx}.
\]
\end{definition}

\begin{theorem}[неравенство Коши--Буняковского--Шварца]
Для любых $f,g\in L_2[a,b]$
\[
  \abs{(f,g)}\le\norm{f}_2\cdot\norm{g}_2,\qquad\text{т.е.}\qquad
  \Bigl(\int_a^b fg\dx\Bigr)^{\!2}\le\int_a^b f^2\dx\cdot\int_a^b g^2\dx .
\]
\end{theorem}

\begin{proof}
При любом $t\in\R$ квадратичный по $t$ многочлен
$\norm{f-tg}^2=\norm f^2-2t(f,g)+t^2\norm g^2\ge0$ неотрицателен, значит его дискриминант
$\le0$: $4(f,g)^2-4\norm f^2\norm g^2\le0$, что и даёт неравенство.
\end{proof}

\begin{remark}[фактор-пространство]
Строго говоря, $\norm f_2=0$ не влечёт $f\equiv0$ (функция может быть отлична от нуля на множестве
меры нуль). Поэтому $L_2$ рассматривают как \emph{фактор-пространство}: функции, различающиеся на
множестве меры нуль, отождествляются. На классах эквивалентности $\norm{\cdot}_2$ — уже настоящая
норма.
\end{remark}

\begin{definition}
Евклидово функциональное пространство называется \emph{гильбертовым}, если оно \emph{полно} по
норме $\norm{\cdot}_2$, т.е. всякая фундаментальная последовательность функций сходится (в
$L_2$) к элементу пространства. Пространство $L_2[a,b]$ гильбертово.
\end{definition}

\section{Ортогональные системы. Тригонометрическая система}

\begin{definition}
Система функций $\{\varphi_n(x)\}_{n=1}^\infty$ называется \emph{ортогональной} в $L_2[a,b]$, если
$(\varphi_n,\varphi_m)=0$ при $n\ne m$ и $\norm{\varphi_n}\ne0$. Если вдобавок
$\norm{\varphi_n}=1$, система \emph{ортонормирована}.
\end{definition}

\begin{example}[тригонометрическая система]
Система $\;1,\ \cos x,\ \sin x,\ \cos2x,\ \sin2x,\ \dots,\ \cos nx,\ \sin nx,\ \dots$
ортогональна в $L_2[-\pi,\pi]$.
\end{example}

\begin{proof}
Достаточно проверить попарную ортогональность. Например, при любых $n,m$
\[
  (\cos nx,\sin mx)=\int_{-\pi}^{\pi}\cos nx\,\sin mx\dx
   =\frac12\int_{-\pi}^{\pi}\bigl(\sin(n+m)x+\sin(m-n)x\bigr)\dx=0
\]
(интеграл нечётной функции по симметричному отрезку). Аналогично $(\cos nx,\cos mx)=0$ и
$(\sin nx,\sin mx)=0$ при $n\ne m$. Нормы:
\[
  \norm1^2=2\pi,\qquad \norm{\cos nx}^2=\int_{-\pi}^{\pi}\cos^2nx\dx=\pi,\qquad
  \norm{\sin nx}^2=\int_{-\pi}^{\pi}\sin^2nx\dx=\pi .
\]
\end{proof}

\section{Ряд Фурье. Коэффициенты Фурье}

\begin{definition}
Пусть $\{\varphi_n\}$ — ортогональная система и $f\in L_2[a,b]$. \emph{Рядом Фурье} функции $f$ по
этой системе называют формальный ряд
\[
  f\sim\sum_{n=1}^\infty c_n\varphi_n(x),\qquad c_n=\frac{(f,\varphi_n)}{\norm{\varphi_n}^2}
  \quad(\text{коэффициенты Фурье}).
\]
\end{definition}

Для тригонометрической системы на $[-\pi,\pi]$ ряд Фурье записывают в виде
\[
  f(x)\sim\frac{a_0}{2}+\sum_{n=1}^\infty\bigl(a_n\cos nx+b_n\sin nx\bigr),
\]
\[
  a_n=\frac1\pi\int_{-\pi}^{\pi}f(x)\cos nx\dx,\quad
  b_n=\frac1\pi\int_{-\pi}^{\pi}f(x)\sin nx\dx,\quad
  a_0=\frac1\pi\int_{-\pi}^{\pi}f(x)\dx .
\]

\begin{theorem}[экстремальное свойство коэффициентов Фурье]
Среди всех многочленов $\sum_{n=1}^N\alpha_n\varphi_n$ среднеквадратичное отклонение
$\delta_N=\norm[\big]{f-\sum_{n=1}^N\alpha_n\varphi_n}^2$ достигает минимума тогда и только тогда,
когда $\alpha_n=c_n$ — коэффициенты Фурье.
\end{theorem}

\begin{proof}
Обозначим $S_N=\sum_{n=1}^N\alpha_n\varphi_n$. В силу ортогональности
\[
  \delta_N=\norm{f}^2-2\sum_{n=1}^N\alpha_n(f,\varphi_n)+\sum_{n=1}^N\alpha_n^2\norm{\varphi_n}^2 .
\]
Подставив $(f,\varphi_n)=c_n\norm{\varphi_n}^2$ и выделив полный квадрат,
\[
  \delta_N=\norm{f}^2-\sum_{n=1}^N c_n^2\norm{\varphi_n}^2
   +\sum_{n=1}^N(\alpha_n-c_n)^2\norm{\varphi_n}^2 .
\]
Последняя сумма $\ge0$ и зависит только от $\alpha_n$; минимум достигается при $\alpha_n=c_n$.
\end{proof}

\begin{theorem}[неравенство Бесселя]
Для любой $f\in L_2[a,b]$
\[
  \sum_{n=1}^\infty c_n^2\norm{\varphi_n}^2\le\norm{f}^2 .
\]
Для тригонометрической системы:
$\displaystyle\frac{a_0^2}{2}+\sum_{n=1}^\infty(a_n^2+b_n^2)\le\frac1\pi\int_{-\pi}^{\pi}f^2\dx$.
\end{theorem}

\begin{proof}
При $\alpha_n=c_n$ из предыдущего доказательства
$0\le\delta_N=\norm{f}^2-\sum_{n=1}^N c_n^2\norm{\varphi_n}^2$, откуда
$\sum_{n=1}^N c_n^2\norm{\varphi_n}^2\le\norm{f}^2$ при всех $N$. Переход $N\to\infty$ даёт
неравенство Бесселя. В частности, $c_n^2\norm{\varphi_n}^2\to0$ — коэффициенты Фурье стремятся к
нулю.
\end{proof}

\section{Замкнутость и полнота. Теорема Стеклова}

\begin{definition}
Ортогональная система $\{\varphi_n\}$ \emph{замкнута}, если для всех $f\in L_2[a,b]$ выполнено
\emph{равенство Парсеваля}
\[
  \norm{f}^2=\sum_{n=1}^\infty c_n^2\norm{\varphi_n}^2
\]
(неравенство Бесселя обращается в равенство, т.е.\ $S_N\to f$ по норме $L_2$).
\end{definition}

\begin{definition}
Ортогональная система \emph{полна} в $L_2[a,b]$, если из $(f,\varphi_n)=0$ для всех $n$ следует
$f\equiv0$ (в смысле $L_2$): ни одна ненулевая функция не ортогональна всем $\varphi_n$.
\end{definition}

\begin{theorem}[Стеклов: эквивалентность полноты и замкнутости]
Ортогональная система в $L_2[a,b]$ замкнута тогда и только тогда, когда она полна.
\end{theorem}

\begin{proof}
$(\Rightarrow)$ Пусть система замкнута и $(f,\varphi_n)=0$ для всех $n$, т.е.\ все $c_n=0$. По
равенству Парсеваля $\norm{f}^2=\sum c_n^2\norm{\varphi_n}^2=0$, значит $f\equiv0$ — система полна.

$(\Leftarrow)$ Пусть система полна. Для $f\in L_2$ рассмотрим остаток
$R_N=f-\sum_{n=1}^N c_n\varphi_n$; по неравенству Бесселя $\sum c_n^2\norm{\varphi_n}^2$ сходится,
поэтому $\{S_N\}$ фундаментальна и (полнота $L_2$) сходится к некоторой $g$. Разность $f-g$
ортогональна всем $\varphi_n$, а в силу полноты системы $f-g\equiv0$, т.е.\ $S_N\to f$. Тогда
$\norm{R_N}^2=\norm{f}^2-\sum_{n=1}^N c_n^2\norm{\varphi_n}^2\to0$ — равенство Парсеваля.
\end{proof}

\begin{example}[ортогонализация: многочлены Лежандра]
В $L_2[-1,1]$ система $1,x,x^2,\dots$ не ортогональна; применяя процесс Грама--Шмидта, получаем
ортогональную систему многочленов Лежандра:
\[
  \psi_0=1,\quad \psi_1=x-\frac{(x,1)}{\norm1^2}\cdot1=x,\quad
  \psi_2=x^2-\frac{(x^2,1)}{\norm1^2}\cdot1=x^2-\frac13,\ \dots
\]
\end{example}

\begin{theorem}
Тригонометрическая система $1,\cos x,\sin x,\dots,\cos nx,\sin nx,\dots$ полна (а значит, и
замкнута) в $L_2[-\pi,\pi]$. Следовательно, ряд Фурье любой $f\in L_2$ сходится к $f$ по норме
$L_2$ и выполняется равенство Парсеваля.
\end{theorem}

\section{Сходимость в точке. Ядро Дирихле}

Перейдём от сходимости «в среднем» ($L_2$) к \emph{поточечной} сходимости ряда Фурье. Пусть $f$
$2\pi$-периодична и интегрируема. Частичная сумма
\[
  S_N(x)=\frac{a_0}{2}+\sum_{n=1}^N(a_n\cos nx+b_n\sin nx).
\]

\begin{definition}
\emph{Ядро Дирихле}
\[
  D_N(t)=\frac12+\sum_{n=1}^N\cos nt=\frac{\sin\bigl((N+\tfrac12)t\bigr)}{2\sin\tfrac t2}.
\]
\end{definition}

\begin{theorem}[интегральное представление]
\[
  S_N(x)=\frac1\pi\int_{-\pi}^{\pi}f(x+t)\,D_N(t)\dt,\qquad
  \frac1\pi\int_{-\pi}^{\pi}D_N(t)\dt=1 .
\]
\end{theorem}

\begin{proof}
Подставим выражения для $a_n,b_n$ в $S_N(x)$ и соберём под одним интегралом по $u$:
\[
  S_N(x)=\frac1\pi\int_{-\pi}^{\pi}f(u)\Bigl[\tfrac12+\sum_{n=1}^N\cos n(u-x)\Bigr]\du
   =\frac1\pi\int_{-\pi}^{\pi}f(u)D_N(u-x)\du .
\]
Замена $t=u-x$ и $2\pi$-периодичность дают первую формулу. Замкнутую форму $D_N$ получаем,
домножив $\tfrac12+\sum\cos nt$ на $2\sin\tfrac t2$ и свернув телескопически
($2\sin\tfrac t2\cos nt=\sin(n+\tfrac12)t-\sin(n-\tfrac12)t$). Нормировка следует из
$\int_{-\pi}^{\pi}\cos nt\dt=0$ при $n\ge1$.
\end{proof}

\begin{lemma}[Римана--Лебега]
Если $f$ интегрируема на $[a,b]$, то
$\displaystyle\int_a^b f(x)\sin(\lambda x)\dx\to0$ и $\displaystyle\int_a^b f(x)\cos(\lambda x)\dx\to0$
при $\lambda\to\infty$. В частности, коэффициенты Фурье $a_n,b_n\to0$.
\end{lemma}

\begin{proof}
Для $f\in C^1$ — интегрированием по частям: $\int f\sin\lambda x\dx
=-\tfrac{f\cos\lambda x}{\lambda}\big|+\tfrac1\lambda\int f'\cos\lambda x\dx=O(1/\lambda)\to0$.
Общий случай — приближением интегрируемой $f$ гладкими функциями в среднем (ступенчатые функции
плотны в $L_1$).
\end{proof}

\begin{theorem}[принцип локализации]
Сходимость ряда Фурье в точке $x$ и его сумма зависят только от поведения $f$ в сколь угодно малой
окрестности точки $x$.
\end{theorem}

\begin{proof}
$S_N(x)-S=\tfrac1\pi\int_{-\pi}^{\pi}\bigl(f(x+t)-S\bigr)D_N(t)\dt$. Вне окрестности $\abs t<\delta$
знаменатель $2\sin\tfrac t2$ отделён от нуля, поэтому
$\dfrac{f(x+t)-S}{2\sin(t/2)}$ интегрируема, и по лемме Римана--Лебега вклад
$\int_{\delta\le\abs t\le\pi}\dots\sin\bigl((N+\tfrac12)t\bigr)\dt\to0$. Остаётся лишь интеграл по
$\abs t<\delta$.
\end{proof}

\section{Признаки сходимости: Дини и Липшиц}

\begin{theorem}[признак Дини]
Пусть существует число $S$ такое, что
$\displaystyle\int_0^{\pi}\frac{\abs{f(x+t)+f(x-t)-2S}}{t}\dt<\infty$.
Тогда ряд Фурье функции $f$ сходится в точке $x$ к $S$.
\end{theorem}

\begin{proof}
В силу чётности $D_N$ запишем
$S_N(x)-S=\tfrac1\pi\int_0^{\pi}\bigl(f(x+t)+f(x-t)-2S\bigr)D_N(t)\dt$. Представим
$D_N(t)=\dfrac{\sin((N+\frac12)t)}{2\sin(t/2)}$ и обозначим
$g(t)=\dfrac{f(x+t)+f(x-t)-2S}{2\sin(t/2)}$. Условие Дини (с учётом $\sin\tfrac t2\sim\tfrac t2$)
гарантирует интегрируемость $g$, и по лемме Римана--Лебега
$\int_0^\pi g(t)\sin\bigl((N+\tfrac12)t\bigr)\dt\to0$, т.е.\ $S_N(x)\to S$.
\end{proof}

\begin{theorem}[признак Липшица/Гёльдера]
Если $f$ непрерывна в точке $x$ и удовлетворяет условию Гёльдера
$\abs{f(x+t)-f(x)}\le L\abs t^{\alpha}$ ($0<\alpha\le1$) в окрестности $x$, то ряд Фурье сходится
в этой точке к $f(x)$. В частности, это верно для дифференцируемых (и кусочно-гладких) функций.
\end{theorem}

\begin{proof}
Берём $S=f(x)$. Тогда $\abs{f(x+t)+f(x-t)-2S}\le 2L\abs t^\alpha$, поэтому
$\dfrac{\abs{f(x+t)+f(x-t)-2S}}{t}\le 2L\,t^{\alpha-1}$ интегрируема в нуле ($\alpha>0$) —
выполнено условие Дини.
\end{proof}

\section{Точки разрыва. Сходимость к полусумме}

\begin{theorem}
Пусть $f$ кусочно-гладкая и в точке $x$ имеет разрыв первого рода с односторонними пределами
$f(x\pm0)$. Тогда ряд Фурье сходится к полусумме
\[
  \lim_{N\to\infty}S_N(x)=\frac{f(x+0)+f(x-0)}{2}.
\]
\end{theorem}

\begin{proof}
Применяем признак Дини с $S=\tfrac12\bigl(f(x+0)+f(x-0)\bigr)$: для кусочно-гладкой $f$ разность
$f(x+t)-f(x+0)=O(t)$ справа и $f(x-t)-f(x-0)=O(t)$ слева, поэтому
$f(x+t)+f(x-t)-2S=O(t)$ и подынтегральное выражение $O(1)$ интегрируемо. Значит $S_N(x)\to S$.
\end{proof}

\begin{example}[формула Лейбница для $\pi$]
Разложим в ряд Фурье на $[-\pi,\pi]$ нечётную функцию $f(x)=\tfrac\pi4\sign x$. Тогда $a_n=0$, а
\[
  b_n=\frac1\pi\int_{-\pi}^{\pi}\frac\pi4\sign x\,\sin nx\dx
   =\frac12\int_0^{\pi}\sin nx\dx\cdot\ldots=\frac{1-(-1)^n}{2n},
\]
т.е.\ $b_{2k}=0$, $b_{2k-1}=\dfrac{1}{2k-1}$. Получаем
$\dfrac\pi4\sign x=\sum_{k=1}^\infty\dfrac{\sin(2k-1)x}{2k-1}$. В точке $x=\tfrac\pi2$ функция
непрерывна и равна $\tfrac\pi4$, а $\sin(2k-1)\tfrac\pi2=(-1)^{k+1}$, откуда
\[
  \frac\pi4=1-\frac13+\frac15-\frac17+\dots \qquad\text{(ряд Лейбница).}
\]
\end{example}

\section{Комплексная форма ряда Фурье}

\begin{definition}
Подставив $\cos nx=\tfrac12(e^{inx}+e^{-inx})$, $\sin nx=\tfrac1{2i}(e^{inx}-e^{-inx})$, получаем
\emph{комплексную форму}
\[
  f(x)\sim\sum_{n=-\infty}^{\infty}c_n e^{inx},\qquad
  c_n=\frac1{2\pi}\int_{-\pi}^{\pi}f(x)e^{-inx}\dx .
\]
\end{definition}

Связь с вещественными коэффициентами:
\[
  c_0=\frac{a_0}{2},\qquad c_n=\frac{a_n-ib_n}{2},\qquad c_{-n}=\frac{a_n+ib_n}{2}\quad(n\ge1).
\]
Система $\{e^{inx}\}_{n\in\Z}$ ортогональна в комплексном $L_2[-\pi,\pi]$ со скалярным
произведением $(f,g)=\int_{-\pi}^{\pi}f\bar g\dx$ и $\norm{e^{inx}}^2=2\pi$. Равенство Парсеваля
принимает вид
\[
  \frac1{2\pi}\int_{-\pi}^{\pi}\abs{f(x)}^2\dx=\sum_{n=-\infty}^{\infty}\abs{c_n}^2 .
\]

\begin{example}
Для $f(x)=e^{ax}$ на $[-\pi,\pi]$:
$\displaystyle c_n=\frac1{2\pi}\int_{-\pi}^{\pi}e^{(a-in)x}\dx
=\frac{(-1)^n\operatorname{sh}(a\pi)}{\pi(a-in)}$, откуда восстанавливается вещественное
разложение через $\operatorname{sh}$ и $\operatorname{ch}$.
\end{example}

\end{document}
